\section{Task × Escala de tiempo: la inteligencia científica no es una puntuación}

Para entender por qué el curso elige sistemas multiagente, primero hay que distinguir dos tipos de capacidad: uno es completar cálculos con extrema velocidad y precisión en tareas con límites claros; el otro es navegar objetivos ambiguos, evidencia contradictoria y experimentos costosos a lo largo de múltiples años. El primero puede generar un «superalocal», mientras que el segundo se acerca más a la inteligencia científica superinteligente que menciona el conferenciante.

\subsection{Lea los ejes de coordenadas antes de leer los nombres de los modelos}

El eje vertical del gráfico asciende desde revisiones bibliográficas, experimentos y análisis de datos hasta tesis doctorales, teorías importantes y rupturas paradigmáticas, representando la complejidad de la tarea y su impacto; el eje horizontal se extiende desde minutos hasta una vida entera, indicando el tiempo necesario para que un científico complete una tarea. Este sistema de coordenadas no es una medición precisa, sino un mapa de diseño del sistema: cuanto más hacia arriba y a la derecha, mayor necesidad de memoria a largo plazo, planificación de recursos, descomposición, validación y colaboración social.

\lecturefigure{slide-004.jpg}{Ubicación de AlphaFold en las coordenadas Task × Escala de tiempo}{Video oficial Stanford CS25 V6 Lecture 07 00:13:10}

\begin{knowledgebox}{Lectura de la figura: AlphaFold comprime una categoría específica de tareas}
AlphaFold reduce el cuello de botella que podría requerir años de experimentación para predecir estructuras de proteínas a un flujo de cálculo más corto, generando así un desplazamiento masivo en el eje horizontal. Sin embargo, resuelve problemas específicos con entradas y salidas bien definidas. El gráfico lo coloca en la parte inferior derecha para enfatizar que un «valor superhumano» puede surgir de la compresión temporal de una tarea estrecha, sin implicar que el sistema ya pueda plantear una agenda de investigación completa.
\end{knowledgebox}

\teachervoice{00:10:17--00:12:45, el docente explica con AlphaFold: un sistema puede tener un impacto revolucionario en una tarea científica específica, pero aún no posee la capacidad de generar hipótesis abiertas, sintetizar conocimientos interdisciplinarios o gestionar proyectos de investigación a largo plazo.}

\subsection{¿Qué cubren los LLMs generales y qué les falta?}

El curso coloca posteriormente Gemini, GPT y Claude cerca del rango de minutos a horas, desde revisiones bibliográficas hasta análisis experimentales preliminares. Su cobertura es más amplia que la de AlphaFold, pero la mayoría se queda en tareas más cortas y de menor riesgo. Este gráfico no afirma que los LLM solo puedan realizar trabajos simples; más bien, cuestiona: ¿cómo extender el conocimiento general y la capacidad de razonamiento lingüístico hacia planes científicos que requieren validación durante días, meses o incluso años?

\lecturefigure{slide-005.jpg}{Estado completo de Task × Escala de tiempo tras añadir LLMs generales}{Video oficial Stanford CS25 V6 Lecture 07 00:15:54}

\begin{knowledgebox}{Lectura de la figura: primero compare el «ancho», luego el «alto»}
AlphaFold es estrecho y profundo: la compresión temporal en tareas específicas es enorme. Los LLMs generales son anchos y superficiales: pueden cubrir muchas tareas de conocimiento, pero a menudo carecen de estado a largo plazo, validación estable y un ciclo cerrado experimental. El objetivo del co-científico de IA no es simplemente arrastrar una burbuja hacia la esquina superior derecha, sino combinar modelos, búsqueda, herramientas, memoria y experimentos humanos para que el sistema mantenga un progreso rastreable bajo mayores plazos y mayor complejidad.
\end{knowledgebox}

\begin{knowledgebox}{Términos clave: tres niveles frecuentemente confundidos}
\begin{tabular}{L{0.16\textwidth}L{0.31\textwidth}L{0.36\textwidth}}
\toprule
Nivel & Problema resuelto & Ejemplo en esta lección \\
\midrule
Task model & Optimiza una tarea específica con entradas y salidas fijas & Predicción de estructura de AlphaFold \\
Foundation model & Proporciona capacidades de conocimiento, lenguaje y razonamiento en tareas amplias & Gemini, GPT, Claude \\
Scientific system & Mantiene objetivos a largo plazo, organiza búsqueda, crítica, herramientas y validación & Sistema multiagente AI co-scientist \\
\bottomrule
\end{tabular}
\end{knowledgebox}

\subsection{De AlphaGo se aprende la organización de la búsqueda, no convertir la ciencia en un tablero}

La lección de AlphaGo radica en que, además de las redes de estrategia fuertes, se puede mejorar continuamente los candidatos mediante autojuego, búsqueda y comparación. La ciencia carece de un tablero fijo, reglas completas y señales inmediatas de victoria o derrota; por ello, no se puede transplantar directamente la Búsqueda de Árbol Monte Carlo (Monte Carlo Tree Search). Lo que sí es transferrible es la idea organizativa de \textbf{hacer presionar a múltiples candidatos entre sí, exponer sus debilidades e iterar su evolución}.

Si una ejecución del sistema genera $N_0$ hipótesis originales, y tras $K$ rondas de comparación y evolución, cada ronda mantiene una proporción $\rho_k$ y expande una proporción $\gamma_k$, la escala aproximada de los candidatos se puede escribir como
$$
N_{k+1}=\rho_kN_k+\gamma_kN_k.
$$
Donde:
\begin{itemize}
  \item $N_k$ es el número de candidatos que aún se rastrean en la ronda $k$;
  \item $\rho_k$ indica la intensidad de retener candidatos de alto valor;
  \item $\gamma_k$ representa la proporción de nuevos candidatos generados a partir de combinaciones, mutaciones o expansiones de candidatos existentes;
  \item La fórmula solo describe el ancho de la búsqueda y no implica que la calidad científica aumente monótonamente con $N_k$.
\end{itemize}

\begin{warningbox}{La ciencia no es un juego con recompensa final}
Los juegos de tablero pueden retroalimentar señales claras de victoria o derrota; la ciencia abierta a menudo tarda años en determinar si una hipótesis se confirma, y «novedad», «viabilidad» e «importancia» pueden estar en conflicto. Introducir un torneo en la ciencia solo proporciona prioridades internas, no convierte el Elo en una ley natural.
\end{warningbox}

\subsection{La inteligencia científica superinteligente es una estructura colaborativa}

En la sección anterior ya se explicó que la ciencia carece de recompensas finales tan claras como los juegos de tablero; por ello, esta sección gira hacia la pregunta sistémica de «quién mantiene la búsqueda a largo plazo». El curso utiliza la visión del cerebro humano junto con múltiples colaboradores para expresar una perspectiva sistémica: la ciencia de alto nivel no es que un módulo único proporcione instantáneamente una respuesta, sino que varias líneas de pensamiento procedan en paralelo, se critiquen mutuamente, utilicen herramientas y memoria, y finalmente los humanos integren los candidatos en experimentos reales. El valor del agente proviene de la división del trabajo y la estructura de retroalimentación, no solo de replicar más instancias idénticas del modelo.

\lecturefigure{slide-006.jpg}{Inteligencia científica expresada como un sistema de colaboración multiagente}{Video oficial Stanford CS25 V6 Lecture 07 00:18:36}

\begin{importantbox}{Cinco condiciones para capacidades a nivel sistémico}
Para pasar de tareas cortas a la ciencia a largo plazo, el sistema necesita al menos: objetivos y estado continuos; exploración paralela y eliminación de duplicados; refutación explícita e incertidumbre; herramientas externas y evidencia; y control humano final sobre experimentos, ética y recursos. Faltar cualquiera de estos elementos significa que aumentar la cantidad de agentes podría solo amplificar la repetición, los errores o la carga administrativa de revisión.
\end{importantbox}

\teachervoice{00:17:30--00:19:28, el docente describe el objetivo como un conjunto de agentes capaces de colaborar a largo plazo, no como un modelo que «sabe todo». Esta formulación prioriza los problemas de arquitectura sobre las capacidades del modelo.}

\subsection{Resumen de la sección}

El gráfico Task × Escala de tiempo separa dos tipos de progreso: los modelos específicos pueden comprimir enormemente una tarea científica, y los modelos generales pueden cubrir más tareas cortas, pero la ciencia abierta requiere un sistema a largo plazo. AlphaGo ofrece inspiración organizativa para la búsqueda, comparación y evolución, no una función de evaluación científica. Por tanto, la inteligencia científica superinteligente debe entenderse como una estructura colaborativa constituida por modelos, agentes, herramientas, memoria, expertos humanos y recursos experimentales.
